\chapter{拓扑学}

拓扑学研究的是在连续变形下保持不变的空间性质。和度量空间相比，拓扑空间不一定预先给出距离函数；它只保留“哪些集合应当被看作开集”这一层结构。因此，拓扑学提供了一种比欧氏空间更抽象、也更灵活的语言，用来讨论连续性、极限、连通性、紧致性以及空间之间的等价关系。

从分析学的角度看，拓扑学可以理解为对“开集方法”的彻底抽象。只要知道开集是什么，就可以重新定义闭集、邻域、收敛、连续映射和紧致性。这样一来，许多原本依赖坐标或距离的定理，就能被放到更一般的空间中重新表述。

\section{从度量空间到拓扑空间}

\subsection{度量空间中的开集}

设 $(X,d)$ 是一个度量空间。对 $x\in X$ 和 $r>0$，开球定义为
\[
    B_r(x)=\{y\in X:d(x,y)<r\}.
\]
子集 $U\subseteq X$ 称为开集，如果对每个 $x\in U$，都存在 $r>0$，使得
\[
    B_r(x)\subseteq U.
\]
这个定义只关心点附近是否能放入一个足够小的开球。换句话说，开集表达的是一种“局部可活动空间”：在 $U$ 中的任一点附近，都有一小块区域仍然完全留在 $U$ 内。

\begin{example}
在通常度量下，区间 $(0,1)$ 是 $\mathbb{R}$ 中的开集；而 $[0,1]$ 不是开集，因为端点 $0,1$ 附近的任意小开区间都会伸出 $[0,1]$ 之外。
\end{example}

\begin{proposition}[度量空间中连续性的拓扑形式]
设 $(X,d_X)$ 和 $(Y,d_Y)$ 是度量空间。映射 $f:X\to Y$ 连续，当且仅当对 $Y$ 中任意开集 $V$，原像
\[
    f^{-1}(V)=\{x\in X:f(x)\in V\}
\]
都是 $X$ 中的开集。
\end{proposition}

这个命题是通往一般拓扑空间的关键。它表明，连续性可以完全用开集的原像来描述，而不必显式使用距离。因此，如果某个集合 $X$ 上已经指定了哪些子集是开集，那么就可以谈论连续性。

\section{拓扑空间}

\subsection{拓扑的定义}

\begin{definition}[拓扑]
设 $X$ 是一个集合。$X$ 上的一个\textbf{拓扑}是 $\mathcal{T}\subseteq\mathcal{P}(X)$，满足：
\begin{enumerate}
    \item $\varnothing\in\mathcal{T}$ 且 $X\in\mathcal{T}$；
    \item 任意多个 $\mathcal{T}$ 中集合的并仍在 $\mathcal{T}$ 中；
    \item 有限多个 $\mathcal{T}$ 中集合的交仍在 $\mathcal{T}$ 中。
\end{enumerate}
二元组 $(X,\mathcal{T})$ 称为\textbf{拓扑空间}。$\mathcal{T}$ 中的元素称为开集。
\end{definition}

这个定义把度量空间中开集的基本性质抽象出来。注意第二条允许任意并，而第三条只要求有限交。原因是任意多个开集的交不一定开，例如
\[
    \bigcap_{n=1}^{\infty}\left(-\frac1n,\frac1n\right)=\{0\}
\]
在 $\mathbb{R}$ 的通常拓扑中就不是开集。

\begin{example}
若 $(X,d)$ 是度量空间，把所有由 $d$ 定义的开集收集起来，就得到 $X$ 上的一个拓扑，称为由度量 $d$ 诱导的拓扑。
\end{example}

\subsection{基本例子}

\begin{example}[离散拓扑]
对任意集合 $X$，令
\[
    \mathcal{T}=\mathcal{P}(X).
\]
也就是说，$X$ 的每个子集都是开集。这个拓扑称为\textbf{离散拓扑}。在离散拓扑中，每个点都是孤立的。
\end{example}

\begin{example}[平凡拓扑]
令
\[
    \mathcal{T}=\{\varnothing,X\}.
\]
这个拓扑称为\textbf{平凡拓扑}或不可分拓扑。它只承认空集和全集是开集，因此区分点的能力极弱。
\end{example}

\begin{example}[$\mathbb{R}^n$ 上的通常拓扑]
$\mathbb{R}^n$ 上由欧氏距离
\[
    d(x,y)=\sqrt{\sum_{i=1}^n(x_i-y_i)^2}
\]
诱导出的拓扑称为通常拓扑。分析学中讨论极限、连续和微分时，默认使用的基本上都是这个拓扑。
\end{example}

\begin{example}[余有限拓扑]
设 $X$ 是无限集合。令 $\mathcal{T}$ 由 $\varnothing$ 以及所有补集有限的集合组成，即
\[
    \mathcal{T}=\{\varnothing\}\cup\{U\subseteq X:X\setminus U\text{ 是有限集}\}.
\]
这给出一个拓扑，称为\textbf{余有限拓扑}。它说明，拓扑空间可以远比度量空间怪异。
\end{example}

\begin{example}[特定点拓扑]
固定一点 $p\in X$。令开集为 $\varnothing$ 以及所有包含 $p$ 的子集。也就是说，
\[
    \mathcal{T}=\{\varnothing\}\cup\{U\subseteq X:p\in U\}.
\]
这是一个拓扑，称为特定点拓扑。
\end{example}

\subsection{拓扑的比较}

\begin{definition}[更细与更粗]
设 $\mathcal{T}_1$ 和 $\mathcal{T}_2$ 是同一集合 $X$ 上的两个拓扑。如果
\[
    \mathcal{T}_1\subseteq\mathcal{T}_2,
\]
则称 $\mathcal{T}_2$ 比 $\mathcal{T}_1$ \textbf{更细}，也称 $\mathcal{T}_1$ 比 $\mathcal{T}_2$ \textbf{更粗}。
\end{definition}

更细的拓扑有更多开集，因此可以表达更强的局部区分能力。离散拓扑是最细的拓扑，平凡拓扑是最粗的拓扑。

\begin{example}
在任意集合 $X$ 上，都有
\[
    \{\varnothing,X\}\subseteq\mathcal{T}\subseteq\mathcal{P}(X),
\]
其中 $\mathcal{T}$ 是 $X$ 上任意拓扑。
\end{example}

\begin{remark}
拓扑越细，映射从该空间出发时越容易连续；拓扑越粗，映射进入该空间时越容易连续。这是因为连续性要求开集原像为开集。
\end{remark}

\section{基与子基}

\subsection{拓扑的基}

在实际使用中，直接列出所有开集通常很困难。更常见的做法是给出一族较小的基本开集，再用它们生成全部开集。

\begin{definition}[基]
设 $X$ 是集合。子集族 $\mathcal{B}\subseteq\mathcal{P}(X)$ 称为 $X$ 上某个拓扑的\textbf{基}，如果：
\begin{enumerate}
    \item 对每个 $x\in X$，存在 $B\in\mathcal{B}$ 使得 $x\in B$；
    \item 若 $x\in B_1\cap B_2$，其中 $B_1,B_2\in\mathcal{B}$，则存在 $B_3\in\mathcal{B}$，使得
    \[
        x\in B_3\subseteq B_1\cap B_2.
    \]
\end{enumerate}
由 $\mathcal{B}$ 生成的拓扑定义为：$U\subseteq X$ 开，当且仅当对每个 $x\in U$，存在 $B\in\mathcal{B}$ 使得 $x\in B\subseteq U$。
\end{definition}

\begin{example}
$\mathbb{R}$ 上所有开区间 $(a,b)$ 构成通常拓扑的一组基。事实上，每个通常意义下的开集都可以写成若干开区间的并。
\end{example}

\begin{proposition}
由一组基 $\mathcal{B}$ 生成的开集族确实构成一个拓扑。
\end{proposition}

\subsection{子基}

\begin{definition}[子基]
设 $X$ 是集合，$\mathcal{S}\subseteq\mathcal{P}(X)$ 且 $\bigcup_{S\in\mathcal{S}}S=X$。由 $\mathcal{S}$ 的有限交作为基生成的拓扑，称为由 $\mathcal{S}$ 生成的拓扑；$\mathcal{S}$ 称为这个拓扑的一组\textbf{子基}。
\end{definition}

\begin{example}
在 $\mathbb{R}$ 上，所有形如 $(-\infty,a)$ 和 $(a,\infty)$ 的集合构成通常拓扑的一组子基。有限交可以给出开区间 $(a,b)$。
\end{example}

\subsection{可数性公理}

\begin{definition}[第一可数与第二可数]
拓扑空间 $X$ 称为\textbf{第一可数}，如果每个点都有一组可数的邻域基。$X$ 称为\textbf{第二可数}，如果它的拓扑有一组可数基。
\end{definition}

\begin{example}
$\mathbb{R}^n$ 是第二可数的。所有以有理点为中心、以正有理数为半径的开球构成一组可数基。
\end{example}

\begin{proposition}
第二可数空间一定是第一可数空间。
\end{proposition}

\section{闭集、闭包、内部与边界}

\begin{definition}[闭集]
设 $(X,\mathcal{T})$ 是拓扑空间。子集 $F\subseteq X$ 称为闭集，如果 $X\setminus F$ 是开集。
\end{definition}

\begin{proposition}[闭集的性质]
闭集满足以下性质：
\begin{enumerate}
    \item $\varnothing$ 和 $X$ 都是闭集；
    \item 任意多个闭集的交仍为闭集；
    \item 有限多个闭集的并仍为闭集。
\end{enumerate}
\end{proposition}

这些性质与开集性质互为对偶，本质上来自德摩根律。

\subsection{闭包与内部}

\begin{definition}[闭包]
设 $A\subseteq X$。包含 $A$ 的所有闭集的交称为 $A$ 的\textbf{闭包}，记作 $\overline{A}$。
\end{definition}

\begin{definition}[内部]
包含在 $A$ 中的所有开集的并称为 $A$ 的\textbf{内部}，记作 $A^\circ$ 或 $\operatorname{Int}(A)$。
\end{definition}

\begin{definition}[边界]
集合 $A$ 的\textbf{边界}定义为
\[
    \partial A=\overline{A}\setminus A^\circ.
\]
等价地，$x\in\partial A$ 当且仅当 $x$ 的每个邻域都同时与 $A$ 和 $X\setminus A$ 相交。
\end{definition}

\begin{example}
在 $\mathbb{R}$ 的通常拓扑中，若 $A=(0,1)$，则
\[
    \overline{A}=[0,1],\qquad A^\circ=(0,1),\qquad \partial A=\{0,1\}.
\]
若 $A=\mathbb{Q}$，则
\[
    \overline{\mathbb{Q}}=\mathbb{R},\qquad \mathbb{Q}^\circ=\varnothing.
\]
\end{example}

\begin{proposition}[用开邻域刻画闭包]
点 $x\in\overline{A}$ 当且仅当 $x$ 的每个开邻域都与 $A$ 相交。
\end{proposition}

\subsection{聚点与稠密性}

\begin{definition}[聚点]
设 $A\subseteq X$。点 $x\in X$ 称为 $A$ 的\textbf{聚点}，如果 $x$ 的每个邻域都含有某个不同于 $x$ 的 $A$ 中点。
\end{definition}

\begin{definition}[稠密子集]
如果 $\overline{A}=X$，则称 $A$ 在 $X$ 中\textbf{稠密}。
\end{definition}

\begin{example}
$\mathbb{Q}$ 和 $\mathbb{R}\setminus\mathbb{Q}$ 都在 $\mathbb{R}$ 中稠密。这个事实说明，稠密性不是“占据大多数点”的意思，而是指每个非空开集都能碰到这个集合。
\end{example}

\section{构造新的空间}

\subsection{子空间拓扑}

\begin{definition}[子空间拓扑]
设 $(X,\mathcal{T})$ 是拓扑空间，$Y\subseteq X$。定义
\[
    \mathcal{T}_Y=\{U\cap Y:U\in\mathcal{T}\}.
\]
则 $\mathcal{T}_Y$ 是 $Y$ 上的拓扑，称为从 $X$ 继承来的\textbf{子空间拓扑}。
\end{definition}

\begin{example}
区间 $[0,1]$ 作为 $\mathbb{R}$ 的子空间时，集合 $[0,1/2)$ 在 $[0,1]$ 中是开集，因为
\[
    [0,1/2)=(-1/2,1/2)\cap[0,1].
\]
但它不是 $\mathbb{R}$ 中的开集。
\end{example}

\subsection{积拓扑}

\begin{definition}[积拓扑]
设 $X,Y$ 是拓扑空间。$X\times Y$ 上的积拓扑以所有
\[
    U\times V,
\]
为基，其中 $U\subseteq X$ 开，$V\subseteq Y$ 开。
\end{definition}

\begin{example}
$\mathbb{R}^2$ 的通常拓扑等于 $\mathbb{R}\times\mathbb{R}$ 的积拓扑。
\end{example}

\newpage

\begin{proposition}[映入积空间的连续性]
映射 $f:Z\to X\times Y$ 连续，当且仅当它的两个分量映射
\[
    \pi_X\circ f:Z\to X,
    \qquad
    \pi_Y\circ f:Z\to Y
\]
都连续。
\end{proposition}

\subsection{商拓扑}

\begin{definition}[商拓扑]
设 $X$ 是拓扑空间，$\sim$ 是 $X$ 上的等价关系，并设 $q:X\to X/\sim$ 是自然投影。集合 $X/\sim$ 上的商拓扑定义为：$U\subseteq X/\sim$ 开，当且仅当
\[
    q^{-1}(U)
\]
在 $X$ 中开。
\end{definition}

\begin{example}[圆作为商空间]
把闭区间 $[0,1]$ 的两个端点粘合，也就是规定 $0\sim1$，得到的商空间同胚于圆 $S^1$。
\end{example}

\begin{remark}
商拓扑形式化了“粘合空间”的想法。圆、环面、莫比乌斯带等空间都可以通过适当的粘合构造出来。
\end{remark}

\section{分离公理}

拓扑空间中的点未必能被开集很好地区分。分离公理用于衡量这种区分能力。

\begin{definition}[$T_0$、$T_1$ 与 Hausdorff 空间]
设 $X$ 是拓扑空间。
\begin{enumerate}
    \item $X$ 称为 $T_0$ 空间，如果任意两个不同点中，至少有一个点有开邻域不包含另一个点；
    \item $X$ 称为 $T_1$ 空间，如果任意两个不同点 $x,y$，都存在开集 $U,V$，使得 $x\in U,y\notin U$ 且 $y\in V,x\notin V$；
    \item $X$ 称为 Hausdorff 空间或 $T_2$ 空间，如果任意两个不同点 $x,y$，存在不相交开集 $U,V$，使得 $x\in U$ 且 $y\in V$。
\end{enumerate}
\end{definition}

\begin{proposition}
Hausdorff 空间一定是 $T_1$ 空间，$T_1$ 空间一定是 $T_0$ 空间。
\end{proposition}

\begin{example}
任意度量空间都是 Hausdorff 空间。若 $x\ne y$，取
\[
    r=\frac12 d(x,y),
\]
则 $B_r(x)$ 和 $B_r(y)$ 不相交。
\end{example}

\begin{proposition}
拓扑空间是 $T_1$ 空间，当且仅当每个单点集 $\{x\}$ 都是闭集。
\end{proposition}

\begin{theorem}[Hausdorff 空间中极限唯一]
在 Hausdorff 空间中，若一个网或序列的极限存在，则极限唯一。特别地，在度量空间中，收敛序列的极限唯一。
\end{theorem}

\section{连续性与同胚}

\subsection{连续映射}

\begin{definition}[连续性]
设 $X,Y$ 是拓扑空间。映射 $f:X\to Y$ 称为连续，如果对任意开集 $V\subseteq Y$，原像 $f^{-1}(V)$ 都是 $X$ 中的开集。
\end{definition}

\begin{proposition}[闭集形式]
映射 $f:X\to Y$ 连续，当且仅当对任意闭集 $F\subseteq Y$，原像 $f^{-1}(F)$ 都是 $X$ 中的闭集。
\end{proposition}

\begin{proposition}[连续映射的复合]
若 $f:X\to Y$ 和 $g:Y\to Z$ 连续，则 $g\circ f:X\to Z$ 连续。
\end{proposition}

\subsection{同胚}

\begin{definition}[同胚]
双射 $f:X\to Y$ 称为\textbf{同胚}，如果 $f$ 连续且反函数 $f^{-1}:Y\to X$ 也连续。若存在同胚 $X\to Y$，则称 $X$ 与 $Y$ 同胚，记作
\[
    X\cong Y.
\]
\end{definition}

同胚是拓扑学中的“相同”。两个同胚空间可能在几何形状上看起来不同，但它们的开集结构完全等价，因此具有相同的拓扑性质。

\begin{example}
开区间 $(0,1)$ 与 $\mathbb{R}$ 同胚。例如映射
\[
    f(x)=\tan\left(\pi x-\frac{\pi}{2}\right)
\]
给出一个同胚 $(0,1)\to\mathbb{R}$。
\end{example}

\begin{remark}
连续双射不一定是同胚，因为反函数可能不连续。因此，在拓扑学中“双射”本身不足以表示两个空间相同。
\end{remark}

\section{连通性}

\subsection{定义}

\begin{definition}[分离]
拓扑空间 $X$ 的一个\textbf{分离}是两个非空、不相交开集 $U,V$，使得
\[
    X=U\cup V.
\]
\end{definition}

\begin{definition}[连通空间]
如果拓扑空间 $X$ 不存在分离，则称 $X$ 是\textbf{连通的}。
\end{definition}

\begin{proposition}[等价刻画]
拓扑空间 $X$ 连通，当且仅当 $X$ 中既开又闭的子集只有 $\varnothing$ 和 $X$ 本身。
\end{proposition}

\begin{theorem}[区间的连通性]
$\mathbb{R}$ 中的任意区间都是连通的。
\end{theorem}

\subsection{连通性的连续像}

\begin{theorem}
连通空间在连续映射下的像仍然连通。也就是说，如果 $f:X\to Y$ 连续且 $X$ 连通，则 $f(X)$ 连通。
\end{theorem}

\begin{theorem}[介值定理]
设 $X$ 连通，$f:X\to\mathbb{R}$ 连续。如果 $a,b\in f(X)$ 且 $a<c<b$，则 $c\in f(X)$。
\end{theorem}

这个定理说明，传统微积分中的介值定理本质上是连通性的结果。

\subsection{道路连通}

\begin{definition}[道路连通]
拓扑空间 $X$ 称为\textbf{道路连通}，如果对任意 $x,y\in X$，都存在连续映射
\[
    \gamma:[0,1]\to X
\]
使得 $\gamma(0)=x$ 且 $\gamma(1)=y$。
\end{definition}

\begin{proposition}
道路连通空间一定连通。
\end{proposition}

\begin{remark}
连通空间不一定道路连通。经典例子是拓扑学家的正弦曲线及其闭包。这说明，连通性只阻止空间被开集拆成两块，而道路连通还要求任意两点之间存在连续路径。
\end{remark}

\section{紧致性}

\subsection{开覆盖}

\begin{definition}[开覆盖]
设 $A\subseteq X$。一族开集 $\{U_\alpha\}_{\alpha\in I}$ 称为 $A$ 的\textbf{开覆盖}，如果
\[
    A\subseteq\bigcup_{\alpha\in I}U_\alpha.
\]
若存在有限多个指标 $\alpha_1,\dots,\alpha_n$ 使得
\[
    A\subseteq U_{\alpha_1}\cup\cdots\cup U_{\alpha_n},
\]
则称这是一组有限子覆盖。
\end{definition}

\begin{definition}[紧致性]
拓扑空间 $X$ 称为\textbf{紧致}，如果 $X$ 的每个开覆盖都有有限子覆盖。子集 $K\subseteq X$ 称为紧致，如果它作为子空间是紧致的。
\end{definition}

\begin{example}
有限拓扑空间总是紧致的，因为任意开覆盖只需要从每个点所在的开集中选出有限多个即可。
\end{example}

\begin{example}
开区间 $(0,1)$ 在通常拓扑下不是紧致的。开覆盖
\[
    \left\{\left(\frac1n,1\right):n\ge2\right\}
\]
覆盖 $(0,1)$，但没有有限子覆盖。
\end{example}

\subsection{$\mathbb{R}^n$ 中的紧致性}

\begin{theorem}[Heine--Borel 定理]
$\mathbb{R}^n$ 中的子集紧致，当且仅当它闭且有界。
\end{theorem}

\begin{remark}
闭且有界推出紧致，是欧氏空间的特殊性质。在一般度量空间或拓扑空间中，闭且有界未必意味着紧致。
\end{remark}

\begin{theorem}[Bolzano--Weierstrass 定理]
$\mathbb{R}^n$ 中任意有界序列都有收敛子列。
\end{theorem}

\begin{remark}
在欧氏空间中，紧致性、序列紧致性以及有界闭集之间有非常紧密的联系；但在一般拓扑空间中，这些概念可能分离。
\end{remark}

\subsection{紧致性的连续像}

\begin{theorem}
紧致空间在连续映射下的像仍然紧致。
\end{theorem}

\begin{theorem}[极值定理]
若 $K$ 紧致，$f:K\to\mathbb{R}$ 连续，则 $f$ 在 $K$ 上取得最大值和最小值。
\end{theorem}

\begin{theorem}[从紧致到 Hausdorff]
若 $X$ 紧致，$Y$ Hausdorff，且 $f:X\to Y$ 是连续双射，则 $f$ 是同胚。
\end{theorem}

\newpage

\section{重新看度量空间}

\begin{definition}[可度量化空间]
拓扑空间 $(X,\mathcal{T})$ 称为\textbf{可度量化}，如果存在度量 $d$，使得 $d$ 诱导出的拓扑恰好等于 $\mathcal{T}$。
\end{definition}

\begin{example}
离散拓扑总是可度量化的。定义
\[
    d(x,y)=
    \begin{cases}
    0,&x=y,\\
    1,&x\ne y,
    \end{cases}
\]
则由该度量诱导的拓扑就是离散拓扑。
\end{example}

\begin{remark}
不是每个拓扑空间都可度量化。度量空间具有许多额外性质，例如 Hausdorff 性、第一可数性等。可度量化定理研究的是哪些拓扑条件足以保证一个拓扑来自某个度量。
\end{remark}

\section{结语}

拓扑学把“接近”“连续”“边界”“整体连通”等观念从坐标和距离中解放出来。它提供了一种高度抽象但非常稳定的语言：只要知道开集，就能重建大量分析学中熟悉的概念。后续的复分析、微分几何、泛函分析乃至现代物理中的空间理论，都离不开这种拓扑语言。

\textbf{关键词：} 拓扑空间，开集，基，闭包，连续映射，同胚，连通性，紧致性。
